\documentclass[11pt]{article}
\usepackage[margin=1in]{geometry}
\usepackage{amsmath,amssymb,amsthm,mathtools}
\usepackage[expansion=false]{microtype}
\usepackage[colorlinks=true,linkcolor=blue,citecolor=blue,urlcolor=blue]{hyperref}
\newtheorem{theorem}{Theorem}[section]
\newtheorem{proposition}[theorem]{Proposition}
\newtheorem{lemma}[theorem]{Lemma}
\theoremstyle{remark}\newtheorem{remark}[theorem]{Remark}
\DeclareMathOperator{\sgn}{sgn}
\DeclareMathOperator{\sig}{sig}
\DeclareMathOperator{\CS}{CS}
\DeclareMathOperator{\Ad}{Ad}
\DeclareMathOperator{\sinc}{sinc}
\newcommand{\R}{\mathbb R}
\newcommand{\Z}{\mathbb Z}
\newcommand{\ii}{\mathrm i}
\newcommand{\E}{\mathbb E}
\newcommand{\abs}[1]{\lvert#1\rvert}
\title{Chern--Simons Phases and Degenerate Sewing\\in \texorpdfstring{$\mathrm{SU}(2)$}{SU(2)} WRT Invariants}
\author{}
\date{}
\begin{document}
\maketitle
\begin{abstract}
We compare exact oriented Witten--Reshetikhin--Turaev sewing with classical flat-connection data in three explicit regimes. In the opposite-cap H plumbing, the closed leading logarithm localizes in the actual quaternionic holonomy chart, where the Coulomb-reduced Chern--Simons action is an exact split cubic. The local density and its scale are obtained from the finite surgery sum. For opposite order-two caps on a graph, a deleted trace computes the absolute adjoint APS invariant, including singular characters; for closed odd cycles the signed leading coefficient is a convergent average of the actual one-loop torsion density, up to its fixed normalization. On a two-seam star, two noncentral branches meet the same central character with opposite APS jumps. Their individual boundary weights diverge, but the signed Weyl sum has a finite limit and differs from the independently computed central coefficient by a factor of \(i\). A separately marked calculation shows how an endpoint determinant changes when the sewing contains a loop. These are regime-specific comparisons, not a general stationary-phase formula for closed WRT invariants.
\end{abstract}

\section{Question and conventions}

A finite torus sewing computes a complex number. Its semiclassical interpretation in terms of flat connections goes back to Witten~\cite{Witten}. A classical comparison should explain its action, its orientation-sensitive phase, and how the comparison changes when a noncentral fibre becomes central. These are different questions. In particular, the RT framing correction is not an adjoint eta invariant, and an equality of actions does not infer a density on a singular character variety. We keep the finite contraction and the classical spectral calculation separate until each comparison is made. Earlier marked Wilson sewing in the Turaev--Viro theory~\cite{Paper42B} motivates this distinction, but none of the closed oriented computations below use a doubled invariant.

Write \(r=k+2\) for shifted level, \(1\le a,b<r\) for charges, and set
\[
 S_{ab}=\sqrt{2/r}\sin(\pi ab/r),\qquad
 \Theta_{aa}=\exp(\pi\ii(a^2-1)/(2r)),\qquad
 \kappa_r=e^{3\pi\ii/4-3\pi\ii/(2r)}.
\]
We normalize \(Z_r(S^2\times S^1)=1\), \(Z_r(S^3)=S_{11}\); in the alternative convention \(\tau_r=Z_r/S_{11}\). The operator defining the real adjoint rho invariant is
\[
 {\cal D}_A(\phi,\alpha)=(d_A^*\alpha,d_A\phi+*_3d_A\alpha),
 \qquad \rho_{\rm ad}=\eta({\cal D}_A)-3\eta({\cal D}_{\rm triv}).
\]
Zero eigenvalues are omitted from eta and from signatures. The Chern--Simons convention uses the fundamental trace and the physical oscillation \(e^{2\pi\ii k\CS}\).

For each graph vertex \(v\) take a punctured sphere times a circle, with one torus for each incident edge and for each listed nonzero cap weight \(p_{vj}\). The spine framing is zero. Coordinates \((z,h)\) on a torus are induced section and fibre, \(z\cdot h=1\). A cap of Kirby weight \(p\) has slope \(-pz+h\); the marked edge \(\epsilon_e S\), \(\epsilon_e=\pm1\), exchanges \(z,h\) with a simultaneous sign \(\epsilon_e\). These data define actual torus gluings even if the graph contains cycles. Let \(A_\epsilon\) be the signed adjacency and put \(Q=\operatorname{diag}(D_v)-A_\epsilon\), \(D_v=-\sum_j1/p_{vj}\).

\begin{lemma}[Exact modular contraction]\label{lem:exact}
Let \(d_v\) be the spine degree and \(c_v\) the number of caps at \(v\). Put \(C_p(a)=(S\Theta^pS)_{a1}\) and \(s=\sum_{v,j}\sgn p_{vj}\). Then the unmarked graph manifold has
\begin{equation}\label{eq:finite}
 Z_r=\kappa_r^{-s-\sig Q}
 \sum_{\boldsymbol a}\prod_v
 \frac{\prod_j C_{p_{vj}}(a_v)}{S_{1a_v}^{\,d_v+c_v-2}}
 \prod_{vw}S_{a_va_w}.
\end{equation}
For any signed nonzero integer \(p\),
\begin{equation}\label{eq:rec}
 C_p(a)=\ii\sqrt{\frac{2}{r|p|}}\,
 e^{\pi\ii\sgn p/4-\pi\ii(p+1/p)/(2r)}
 \sum_{j\bmod|p|}
    \sin\frac{\pi(a+2rj)}{rp}
    e^{-\pi\ii(a+2rj)^2/(2rp)} .
\end{equation}
\end{lemma}
\begin{proof}
Verlinde sewing gives the vertex denominator and the cap column. With the section Lagrangian on all ports, weight-zero cylinders and caps have matrices \(S\) and columns \(\Theta^p S|1\rangle\) before the incidence \(S\). One verifies these weights by testing against section-meridian solid tori: the test kernel equals the declared Lagrangian, so its Maslov correction vanishes. Signed swaps give the same \(S\) because all labels are self-dual. The actual simultaneous-interface seam kernel has the form \((T y,y)\), where a cap contributes the diagonal entry \(-p\) and an edge the hyperbolic block with off-diagonal \(-\epsilon_e\). If \(H\) sums incidence coordinates at each vertex, the product-block kernel is \(\{(H^{\mathsf T}x,y):Hy=0\}\). The interface Maslov form is \(T|_{\ker H}\), whose signature equals \(\sig T-\sig(HT^{-1}H^{\mathsf T})=-s-\sig Q\). The inertia identity holds also in the singular case: the common radical of the two complementary restrictions splits off in a hyperbolic summand. The unextended gluing rule multiplies by \(\kappa_r\) to this signature~\cite[IV.3.5, IV.7.1]{Turaev}; see also~\cite{GilmerMasbaum}. For \eqref{eq:rec}, the two sine products turn the defining sum into \(e^{-\pi\ii p/(2r)}[G(a-1)-G(a+1)]/(2r)\), where \(G(u)=\sum_{b\bmod 2r}e^{\pi\ii p b^2/(2r)+\pi\ii ub/r}\). Apply signed quadratic reciprocity~\cite[Prop.\ 2.3]{Jeffrey}; subtracting the neighboring squares gives \eqref{eq:rec}.
\end{proof}

\section{An exact split cubic in closed H sewing}\label{sec:H}

Let \(M_H\) be the six-component H plumbing: join two zero-framed centres by one Hopf edge; at each centre attach the two meridian leaves of framings \(+2,-2\). All nonzero linking magnitudes are one, and no other components link. We use positive Hopf convention for the spine, equivalent to a signed exchange of the induced peripheral axes. The final scalar here is independent of this sign. Denote by \(\rho_H\) the quaternionic character described below.

\begin{theorem}[Localized closed logarithm]\label{thm:H}
At \(\rho_H\), the harmonic Coulomb-reduced Chern--Simons action is exactly a nonzero split real cubic \(cXYZ\), and \(\CS(\rho_H)=1/2\pmod1\). In arbitrarily small fixed neighborhoods of this chart the original unmarked closed level sum has the representation
\begin{equation}\label{eq:localH}
 Z_r(M_H)=(-1)^r r^{3/2}\!
   \int_{\R^3} e^{2\pi\ii r cXYZ}\,\varrho(X,Y,Z)\,dX\,dY\,dZ
     +O(\sqrt r),
 \qquad \varrho(0)=\frac{\sqrt2\,\operatorname{Vol}(M_H)}{2\pi^3}.
\end{equation}
Here \(\varrho\) is a smooth local density derived from the surgery expression with fixed smooth cutoffs. One may take \(|c|=\operatorname{Vol}(M_H)/\pi^2\) in orthonormal parallel coordinates. Consequently the entire leading term of the closed scalar is \((-1)^r\sqrt{2r}\log r/\pi\), or \((-1)^r r^2\log r/\pi^2\) for \(\tau_r\). Formula \eqref{eq:localH} does not identify the \(O(\sqrt r)\) part.
\end{theorem}

\subsection{Classical chart and the action value}

Each half of the H graph has generating sections \(x,y,z\), fibre \(h\), and relations \(x^2 h=1,\ y^{-2}h=1,\ z=(xy)^{-1}\), after choices of cone meridian signs. There is a marked flat model
\[
 (\R^2\times[-1,1])/\langle A,B\rangle,\quad
 A(X_0,Y_0,t)=(X_0+1,-Y_0,-t),\quad B(X_0,Y_0,t)=(X_0,Y_0+2,t).
\]
The circle fibration in the \(X_0\) direction has disk base with two order-two cone points. The two cone lifts can be chosen as \(x=A,\ y=A^{-1}B\) with \(h=A^{-2}\); indeed their squares give exactly the marked filling slopes, and the product section is \(z=B^{-1}\). The two lifts can be prescribed independently because a section change over the pair of pants has two independent integral circle degrees. This checks the actual peripheral marking, not just its abstract group. On the boundary the exchanged directions have equal periods and product collars, so the marked exchange is a flat isometry. This is the flat torus semi-bundle case~\cite{SWW}. Writing \(C\) for the other glide, van Kampen gives \(AC^2A^{-1}=C^{-2},\ CA^2C^{-1}=A^{-2}\). Sending \(A,C\) to orthogonal imaginary unit quaternions \(\mathbf i,\mathbf j\) gives a representation. Their adjoint half-turns are exactly the tangent linear \(V_4\) representation.

Here is why its cubic is exact. For a flat connection with \(*[\mathcal H^1,\mathcal H^1]\subset\mathcal H^1\), split the Sobolev Coulomb slice as \(a+b\) with \(a\) harmonic and \(b\) coexact. The projected transverse stationarity equation is
\[
 P_{\perp}*_3\big(d_A b+\tfrac12[a+b,a+b]\big)=0 .
\]
Its coexact linearization \(*d_A\) is an isomorphism from order \(s\) to order \(s-1\), \(s>3/2\): its square is \(d_A^*d_A\) on this summand, with no harmonic kernel. The implicit function theorem gives a unique small \(b(a)\); the harmonic bracket condition makes \(b=0\) an exact solution. On a flat compact manifold with flat unitary coefficients, Bochner's identity makes twisted harmonic forms parallel, hence the bracket condition holds. Harmonic one-forms are equivariant linear maps between the tangent and adjoint representations. Both \(V_4\) representations split into the same three distinct real nontrivial characters, so the harmonic form is \(X\mathbf i\,dx_1+Y\mathbf j\,dx_2+Z\mathbf k\,dx_3\). The cubic trace has \(\operatorname{Tr}(\mathbf i\mathbf j\mathbf k)=-2\), giving \(|c|=\operatorname{Vol}/\pi^2\); there is no quadratic term. This representation has no infinitesimal stabilizer. The exact critical set is the three axes. The argument identifies a reduced action, not a measure on all connections.

For the action value, rotate the second generator from \(\mathbf j\) to \(\mathbf i\) with both squares fixed at \(-1\), an actual flat path. The endpoint eigenline sends both plus-cap meridians to \(1/4\in\R/\Z\). The surgery intersection matrix, on the two centres followed by the two pairs of leaves, is
\[
 \begin{pmatrix}0&1&1&1&0&0\\1&0&0&0&1&1\\1&0&2&0&0&0\\
 1&0&0&-2&0&0\\0&1&0&0&2&0\\0&1&0&0&0&-2\end{pmatrix}.
\]
Elimination gives \(H_1=(\Z/4)^2\) generated by the two plus meridians, with linking matrix \(\bigl(\begin{smallmatrix}1/2&1/4\\1/4&1/2\end{smallmatrix}\bigr)\) up to overall orientation. The corresponding linking dual has square \(1/2\) modulo one. Since \(c_2(L\oplus L^{-1})=-c_1(L)^2\), the fundamental-trace transgression of this eigenline gives \(\CS=1/2\pmod1\), unaffected by an overall sign; see~\cite{Auckly}. Constancy along the flat path gives the irreducible action value.

\subsection{Finite localization and density normalization}

From \eqref{eq:rec},
\[
 |C_2(r-n)|^2=\tfrac4r{\bf1}_{n\ {\rm odd}}\cos^2(\pi n/(2r)).
\]
The exact H scalar is thus
\[
 Z_r=(-1)^r(2/r)^{3/2}\sum_{n,m\ {\rm odd}<r}
      \cot(\pi n/(2r))\cot(\pi m/(2r))\sin(\pi nm/r).
\]
For reference the finite sine transform
\[
 \tfrac2r\sum_{n\ {\rm odd}<r}\cot(\pi n/(2r))\sin(\pi nm/r)=1,\qquad 1\le m<r
\]
follows by doubling to the full odd grid modulo \(2r\) and using
\(\sin(mq)\cot(q/2)=1+\cos(mq)+2\sum_{j=1}^{m-1}\cos(jq)\). In particular the exact scalar also equals \((-1)^r\sqrt{2/r}\sum_{m\ {\rm odd}<r}\cot(\pi m/(2r))\). The orbital expression is stronger for localization. Let \(x=n/r,y=m/r,\ h=2/r\), \(a(x)=x\cot(\pi x/2)\), \(a(0)=2/\pi\). Since the relative dot product \(t\) of two Haar-transported adjoint axes is uniform on \([-1,1]\), the edge sine becomes an exact orbital integral, giving
\begin{equation}\label{eq:Horbit}
 Z_r=(-1)^r\frac{\pi\sqrt2}{4}r^{3/2}
 h^2\!\sum_{x,y\ {\rm positive\ odd\ mesh}<1}\!
 a(x)a(y)\int_{-1}^{1}e^{\pi\ii r xyt}\,dt.
\end{equation}

Choose even smooth cutoffs \(\chi(x)\chi(y)\eta(t)\), all one near zero and supported in arbitrarily small fixed intervals. The part omitted from \eqref{eq:Horbit} is \(O(\sqrt r)\). For the exterior charge part, the exact sine transform is the odd \(2r\)-periodic square wave, with zero at its breaks. Multiplication by an even smooth period-two weight \(f(n/r)\) with \(f(0)=0\) convolves it with rapidly decreasing Fourier coefficients. Away from a break the constant terms sum to zero, so the new transform at \(m\) is \(O_N((1+\min(m,r-m))^{-N})\). The remaining cotangent, bounded by \(O(r/m)\) at one end and \(O((r-m)/r)\) at the other, gives the claimed error. This works whenever one cutoff removes its near-zero charges. After both colour cutoffs, signed extension of the two odd grids and even \(t\)-integration converts their positive quadrant to one quarter of the full signed grid. Two-dimensional Poisson summation has phases \(\pi r(txy-u x-vy)\), \(u,v\in\Z\). All nonzero aliases are uniformly nonstationary since the colour support has \(|x|,|y|<d<1\) and \(|t|\le1\); repeated integration by parts sums their errors to \(O_N(r^{-N})\) before the prefactor. The zero alias with \(t\) away from zero is \(O(r^{-1})\) by integrating in \(x\) and using the Schwartz transform in \(ty\). This also removes the angular complement at \(O(\sqrt r)\). The remaining integral is
\[
 (-1)^r\,\frac{\pi\sqrt2}{16}r^{3/2}
 \int_{\R^3} a(x)\chi(x)a(y)\chi(y)\eta(t)e^{\pi\ii r xyt}\,dx\,dy\,dt .
\]

To calibrate its variables, let the spatial glide translations be \((a_0,0,0)\) for \(A\) and \((0,b_0,d_0)\) for \(C\) after choosing a Euclidean origin on the first axis. The squares of \(A,C,AC\) generate the translation lattice: they span the respective independent directions, and the quotient generated by the two resulting involutions with involutive product has order four, the already known linear holonomy order. Thus \(\operatorname{Vol}=2|a_0b_0d_0|\). On the parallel chart the fixed straight representative paths yield holonomies
\[
 U=\mathbf i e^{a_0X\mathbf i},\qquad V=\mathbf j e^{b_0Y\mathbf j+d_0Z\mathbf k},
\]
up to a simultaneous choice of connection signs. Their regular fibres are \(U^{-2},V^{-2}\). The Wilson eigenvalue \(2\cos(\pi a/r)\) identifies the two signed conjugacy-defect lifts with \(x,y\). The third coordinate is the relative dot product \(t\) of the normalized axes of \(U,V\), calibrated by the exact orbital integral. Direct quaternion multiplication gives linear terms \(2a_0X/\pi,\,2b_0Y/\pi,\,d_0Z\) up to signs, and each of \(x,y,t\) vanishes exactly on its corresponding coordinate plane locally. Hence \(2\pi cXYZ=\pi xyt\,g(x,y,t)\) with \(g(0)=\pm1\) by the volume identity. Absorb \(g\) in the local \(t\) coordinate; the Jacobian at zero has absolute value one for this last change. The pushed density at zero is
\(\frac{\pi\sqrt2}{16}(2/\pi)^2\, 4|a_0b_0d_0|/\pi^2\), as in \eqref{eq:localH}. The phase \(r\) versus \(k=r-2\) at \(\CS=1/2\) is the same; the two actions in the local integral differ only by a smooth density multiplier with value one at zero.

Finally for any smooth compactly supported \(g_0\), the integral of \(e^{\ii\lambda cxyz}g_0\) has leading coefficient \(4\pi g_0(0)\log\lambda/(\lambda|c|)\). To see the remainder, subtract an even product cutoff times \(g_0(0)\); its error is a sum of coordinate multiples and has integral \(O(\lambda^{-1})\) by the Fourier bound \(C_N(1+\lambda|xy|)^{-N}\) in the third variable. The product term on a positive square yields \((\lambda|c|)^{-1}\int_0^{\lambda|c|d^2}\widehat\eta(s)\log(\lambda|c|d^2/s)\,ds\), and the four quadrants give the result since \(\int_0^\infty\widehat\eta(s)\,ds=\pi\eta(0)\). Substituting \(\lambda=2\pi r\) proves Theorem~\ref{thm:H}. The angular variable in this proof is an exact orbital variable, not a claim of a finite operator projecting onto arbitrary relative angles.

\section{Absolute phase and the regular sewing density}\label{sec:APS}

Let all spine fibres of the boundary local system \(E\) act trivially, so the same is true of all spine sections. The system can be the rank-three adjoint of a flat \(SU(2)\) connection with central spine fibres. Write \(I_v\) for the invariant subspace of the entire local system at \(v\), and \(R_e\) for its orthogonal or unitary edge transport. Put \(Q_I\) on \(\bigoplus_v I_v\) with diagonal \(D_v I\) and edge pairing \(-\epsilon_e\langle u_v,R_eu_w\rangle\). For complex rank-one systems use \(\rho_E=\eta({\cal D}_E)-\operatorname{rank}(E)\eta({\cal D}_{\rm triv})\).

\begin{proposition}[Deleted-trace formula]\label{prop:aps}
For the oriented central local system above,
\begin{equation}\label{eq:deletion}
 \rho_E(M) = \operatorname{rank}(E)\sig Q-\sig Q_I+
             \sum_{\text{caps}}\rho_{E|_{\partial N}}(\partial N),
\end{equation}
where \(N\) is the actual corresponding terminal disk bundle with its trace orientation. This is an identity of real numbers also with nullities.
For the opposite order-two pair on every vertex, on the branch with all fibres \(-1\), let \(u_v\) be the unit cone axis and \(B_{vw}=\epsilon_e\langle u_v,R_eu_w\rangle\) across an edge. Then
\begin{equation}\label{eq:pairedaps}
 \CS=l/2\pmod1,\qquad \rho_{\rm ad}=\sig B-3\sig A_\epsilon,\qquad l=|E(G)|.
\end{equation}
\end{proposition}

\begin{proof}
Use the actual disk-bundle plumbing trace, with terminal weights \(p\) and edge intersections \(-\epsilon_e\) in the stated peripheral convention. The latter sign follows by exchanging the fibre and base disk coordinates: the section bounding the deleted base disk has the opposite orientation from its base complex angle. Delete the interiors of the small terminal disk bundles \(N\). Locally the surviving central bundle retracts to the punctured sphere; the remaining terminal region retracts to \(\partial N\) and overlaps the central region in a meridian circle up to homotopy. Ordinary graph-edge overlaps are contractible and no two such overlaps meet. As the boundary fibre is trivial in \(E\), the cone representation extends across the local trace by the order-\(|p|\) lens relation. Graph stable transports extend as well, since the edge overlaps are contractible. In the negative-central order-two case the fundamental representation need not extend; the adjoint does extend, with each order-two cone generator a half-turn.

The terminal lens pieces have no degree-one or degree-two homology with these coefficients, by finite transfer over characteristic zero. At a vertex the local Mayer--Vietoris group in degree two is therefore the kernel of the peripheral-circle maps into the punctured base, equivalently \(H_2(P_v,\partial P_v;E)\cong I_v\) by twisted duality (use \(H_2(S^2;E)\) directly when no cap is removed). For an invariant coefficient the representative is the central sphere with its rational terminal corrections \(S_v-\sum_jS_{vj}/p_{vj}\). Its square is \(D_v\); at an edge it uses the unchanged central chart. Intersecting these representatives pairs their transported coefficients with sign \(-\epsilon_e\). Thus the ordinary and twisted middle signatures of the deleted trace are exactly \(\sig Q,\sig Q_I\), with radicals omitted, including when the graph has cycles. APS subtraction in the positive mixed-operator convention gives \(\rho_E(\partial W^\circ)=\operatorname{rank}(E)\operatorname{Sign} W^\circ-\operatorname{Sign}_E W^\circ\)~\cite{APS,KKR}. The extra boundaries are \(-\partial N\), proving \eqref{eq:deletion}. For opposite order-two caps \(I_v=\R u_v\), hence \(Q_I=-B\); the two lens systems are isomorphic with opposite orientations. This proves the rho equality.

For the action on this branch, deform the stable transports to a common axial frame; the cap and central seam relations remain flat. There is an abelian endpoint with spine eigenline coordinates \(-1/2\), plus-leaf coordinates \(1/4\), and minus-leaf coordinates \(s_v/2-1/4\), where \(s_v=\sum_{w\sim v}\epsilon_{vw}\). Let this vector be \(\alpha\), and \(L\) the integral trace matrix (cap incidence can be made \(+1\)). Then \(L\alpha\) is integral: its entries are \(s_v\) on the spine, zero on the plus leaves, and \(-s_v\) on the minus leaves. Thus the eigenline class has the rational relative lift \(\alpha\); its square is \(\alpha^{\mathsf T}L\alpha\) modulo one, valid even if the matrix is singular because the lift exists. The free graph phases can be extended on the trace and do not affect the square. The spine contribution \(-\sum_v s_v/2\) is integral. The remaining term, up to overall sign, is \(\sum_v s_v(2s_v-1)/4\equiv l/2\pmod1\), since \(s_v(s_v-1)\) is even and \(\sum s_v=2\sum_e\epsilon_e\). Eigenline transgression as in Section~\ref{sec:H} and flat continuity give the stated CS.
\end{proof}

\begin{theorem}[Actual one-loop comparison for closed odd cycles]\label{thm:cycle}
For opposite order-two caps and a closed odd spine cycle of length \(n\), let \(\sigma=\sig A_\epsilon\). The all-negative-fibre locus is smooth at its irreducible points with \(\det B\ne0\). There the adjoint Reidemeister half-density is \(c_G|\det B|^{-1/2}d\mu\), with \(c_G>0\) constant and \(d\mu\) quotient Haar. The full unmarked quantum asymptotic, normalized by \(r^n\), is
\begin{equation}\label{eq:cycleaps}
 e^{2\pi\ii k\CS}\frac{2^{n/2}}{\pi^n}\,2^{-n/2}
   \E_\mu\left[\frac{e^{-\pi\ii\rho_{\rm ad}/4}}{\sqrt{|\det B|}}\right]+o(1).
\end{equation}
The expectation is absolutely convergent. In particular it gives \(Z_r/r^n=(-1)^{rn}e^{3\pi\ii\sigma/4}2^{n-1}/\pi^n+o(1)\).
\end{theorem}

\begin{proof}
At each vertex the two cone quaternions are \(u_v\) and a specified \(\pm u_v\); all marked torus holonomies are \(-1\). After fixing the axes, the edge conjugators are free, with one axial \(U(1)\) gauge action per vertex. At a regular point there are no additional nearby fibre defects. Indeed the centrality and two cap square relations require the axial projection of the logarithm of the sign-stripped product of sections to vanish. Its derivative in the fibre angle defects is \(B\); at zero defect these equations hold for all nearby axis and transport choices. The implicit function theorem and the existing flat all-central family therefore identify the local locus before the axial quotient. At irreducible points the axial action has only the finite central stabilizer (trivial in the effective action), so the quotient is smooth there. No such smoothness is claimed at reducibles: aligned transports may have \(\det B\ne0\) and still have an axial stabilizer. Reducible transports have \(T_e=\pm1\) on each edge, since all local cone axes must lie on one transported invariant line, and are consequently Haar-null.

For the torsion density, right-trivialized cocycles give \(H^0\) the axis line at each block. Their boundary \(H^1\) traces are a common fibre variation \(p_vu_v\) and section variations \(s_{ve}\) subject to \(\sum_e\langle u_v,s_{ve}\rangle=0\): the two square relations determine longitudinal cap variations, and their product determines the second transverse variation after killing the first by a coboundary. This restriction is injective and has rank \(3d_v\). \(H^2\) injects into boundary area-coordinates with the same summed axial condition by the pair sequence and duality. Put \(P\) equal to the space of all incidence section vectors and \(Ua=(a_vu_v)_{ve}\). The degree-one gluing map has the form \((s,p)\mapsto s-Kp\) on \(U^{\mathsf T}s=0\), where \(Kp\) is the swapped transported fibre variation. Augmentation by \(U^{\mathsf T}s\) on all of \(P\) leaves determinant \(|\det U^{\mathsf T}K|=|\det B|\). These volumes are fixed block quotient volumes. In degree zero the seam map is \(J={\cal A}U\), where \({\cal A}s\) is the difference of transported incidence vectors. It is injective at an irreducible. In degree two the seam map is the corresponding sum \({\cal C}\) on \(\ker U^{\mathsf T}\); \(\ker{\cal C}=\operatorname{im}{\cal A}^{\mathsf T}\) and \(U^{\mathsf T}{\cal A}^{\mathsf T}=J^{\mathsf T}\). For tangent lifts \(y\) to the seam space, put \(c=|\det[J,y]|\). The Poincaré-dual degree-two basis is \({\cal A}^{\mathsf T}b\), where \(b\in\ker J^{\mathsf T}\) is dual to \(y\); the duality follows by integrating the \(H^1\) connecting representative along the seam collar. Appending to it a corrected one-sided lift of the degree-two seam map gives determinant \(c^{-1}\) by the dual-basis identity. Thus multiplicativity of cochain torsion with dual bases gives \(c^2/|\det B|\), up to constant block determinants (an isolated isomorphism from degree one to degree two contributes the inverse determinant). See~\cite[Prop.\ 3.3]{Massuyeau}. The quotient Haar volume in tangent coordinates is proportional to \(c\), by the axial gauge columns \(J\). Taking a square root proves the density assertion; no absolute normalization of a functional measure is implied.

Independently, let \(d=r-a\). The exact cap identity of Section~\ref{sec:H} and the edge reflection sign on the supported odd defects give
\[
 Z_r=(-1)^{rn}\kappa_r^\sigma\frac{2^{3n/2}}{\pi^n}r^{n/2}
 \sum_{\boldsymbol d\ {\rm odd}<r}\,
  \prod_v \left[\cos^2(\pi d_v/(2r))
             \left(\frac{\pi d_v/r}{\sin(\pi d_v/r)}\right)^2\right]
  \prod_{vw}\sinc(\pi d_vd_w/r).
\]
The bracket is bounded on the entire interval including the upper end. At \(d_v=\sqrt r\,t_v\), each odd lattice gives density \(\sqrt r/2\), yielding coefficient \(2^{n/2}\pi^{-n}{\cal J}\) with \({\cal J}=\int_{t>0}\prod_{vw}\sinc(\pi t_vt_w)\,d^nt\). Uniform tails follow on dyadic boxes \(t_v\asymp2^{j_v}\): the volume times sinc bound is at most \(C2^{-\sum_{vw}|j_v+j_w|/2}\), since the graph has degree two. The edge-sum linear map is invertible on an odd cycle, so these bounds sum at all scales. This is a whole closed limit, not an inserted corner.

The normalized Haar marks \(T_e=\langle u_v,R_eu_w\rangle\) are independent uniform on \([-1,1]\), also around graph cycles: the unconstrained edge conjugators are product Haar before dividing by the constant gauge volume. Write each sinc as this orbital integral; with positive Gaussian damping Fubini holds and the sign-symmetric edge law extends the orthant with factor \(2^{-n}\). Its Gaussian gives \({\cal J}=2^{-n/2}\E_\mu e^{-\pi\ii\sig B/4}/\sqrt{|\det B|}\), choosing the minus sign by the symmetry \(B\mapsto-B\) before the limit. The determinant envelope is integrable: \(\det B=2\prod_e\epsilon_eT_e\), only the two full directed cycles contribute, and \(\E|T_e|^{-1/2}=2\). The signatures are the opposite values \(+1,-1\) on the two sign classes (deform magnitudes to one and diagonalize the cycle), so \({\cal J}=2^{n/2-1}\). Finally combine \(\kappa_r^\sigma\) with \eqref{eq:pairedaps} and the independent action \(e^{2\pi\ii k\CS}=(-1)^{rn}\). This proves \eqref{eq:cycleaps}. The rho identity holds beyond this odd-cycle example but does not itself show integrability or a nonzero leading coefficient on other graphs.
\end{proof}

\section{A signed Weyl collision at a central fibre}\label{sec:wall}

In a zero-Euler triple, write \(t=a/r\) for the unreflected fibre angle in units of \(\pi\), and \(u_j\) for the even reciprocal residues of \eqref{eq:rec}. Thus \(\sum_j1/p_j=0\) and in a resolved bulk term the local resonance is \(\Lambda=\sum_j u_j/p_j\in\Z\), with parity supplied by adjacent corners. Define
\[
 P_u(t)=\prod_{j=1}^3\sin[\pi(u_j+t)/p_j].
\]
At a central wall \(e=0\) or \(1\) the signed Weyl involution \(u\mapsto u'=-u-2e\) preserves the resonance parity. In the endpoint class the axes reverse and \(u_j+e\) becomes its negative. The quadratic frequency changes by \(4e\Lambda\), since the reciprocal sum is zero, so its \(r\)-phase is unchanged. In particular
\begin{equation}\label{eq:weyl-pair}
 P_{u'}(e+s)=-P_u(e-s),\qquad
 \lim_{t\to e}\frac{P_u(t)+P_{u'}(t)}{\sin\pi t}
  =\frac{2(-1)^e}{\pi}P_u'(e).
\end{equation}
This cancels an individual simple pole after signed grouping, not after replacing each residue by its absolute value. We next compare a finite wall value with a distinct physical central charge coefficient.

\begin{theorem}[Closed two-seam star]\label{thm:star}
Take a three-vertex path with positive swaps, central caps \((4,4,-2)\) and four \(+2\) caps on each leaf. There are two actual resolved noncentral branches meeting the same central character at \(t=1\), all with \(\CS=0\). In the orientations above their adjoint rho invariants are respectively \(-3,-7\) in the interior and \(-5\) at the common endpoint. The two branch densities diverge individually; their signed sum has limit \(i\) times the nonzero central coefficient at that tuple. In the full unmarked contraction
\begin{equation}\label{eq:star}
 \frac{Z_r}{r^5}=(1+(-1)^r)\left[
       \frac{e^{5\pi\ii/4}}{4\sqrt2\,\pi^2}
      +\frac{3\zeta(3)e^{-\pi\ii/4}}{2\sqrt2\,\pi^4}\right]+o(1).
\end{equation}
The coefficient on odd levels is zero; no conclusion at lower orders is asserted.
\end{theorem}

\subsection{Actual flat phases}

The leaves have negative fibre and four trace-zero cap holonomies. Use a fixed axis \(\mathbf i\) for the central block, with fibre angle \(t\), and the two residue triples \(u=(0,0,0)\) and \((-2,-2,-2)\). They have \(\Lambda=0\). At \(t=1\), using opposite axes, their cap holonomies become the same signed matrices; the three angles then are \((\pi/4,\pi/4,\pi/2)\). One can keep the ordered four-cap leaf products flat as \(t\) changes: choose the first pair product a pure unit quaternion perpendicular to the central axis, and the other pair product its inverse times the required inverse section. Both products are noncentral, and each admits a circle of trace-zero pair realizations, continuously along these paths. This gives the same endpoint tuple on both sides by the free pair choices. At the central endpoint these four-cap leaf products equal \(-1\). One may vary a first pair product continuously to \(+1\) and the other to \(-1\) along one axis, then align all axes; this is a flat deformation to an abelian endpoint. The central eigenline has spine coordinates \(-1/2\) and cap coordinates \(s_j/(2p_j)\) with \(s_j\) odd, taken as signed unit labels here. The marked trace matrix times this vector is integral by the cap and product relations. Its square modulo one is \(-2/2+\sum s_j^2/(4p_j)\): the omitted sum of cap eigenphases is integral by the product relations (there are two edges). The centre cap contribution is zero, and the two leaf contributions are \(1/2\) each. The square is therefore zero modulo one. By the eigenline transgression and the actual flat paths, the common CS is zero.

The absolute spectral computation can be made without using the quantum phases. At the central endpoint, \eqref{eq:deletion} gives \(3\sig Q=-3\), \(Q_I=0\) (only a central axis with Euler weight zero), and cap lens sum \(-2\). In detail the nontrivial signed order-\(|p|\) adjoint-plane lens correction for residue \(a\) is
\[
 L(p,a)=4a^2/p-4\sgn(p)a+2\sgn(p),\qquad 1\le a<|p|.
\]
A trivial character has zero correction, not the polynomial value at zero. This is the two conjugate lines of the lens cotangent formula~\cite[Sec.\ 2.3]{Toffoli}, in the oriented terminal trace convention of \eqref{eq:deletion}. Here \(L(4,1)=-1\), \(L(\pm2,1)=0\); thus the central value is \(-5\).

For the interior choose \(t=1/2\) and compress both tori at slopes \((1,2)\) viewed from the middle and \((2,1)\) from the leaves. These slopes have trivial fundamental holonomy. To see the absolute defect, use the common-boundary cut-and-paste formula~\cite[Thm.\ 3.9]{Toffoli}. The replaced seam cylinders closed by matching meridians are \(S^2\times S^1\)'s with rho zero (flat extension over \(D^3\times S^1\)). With \(\beta=b/a\) for each filling slope and \(K_v=D_v+\sum\beta\), the component kernel is
\[
 x_{ve}=J_{ve}c_v,\qquad \sum_eJ_{ve}^*y_{ve}=-D_v c_v,
\]
where \(W_e\) is the fixed space of the peripheral local system, \(I_v\) the local fixed space and \(J_{ve}:I_v\to W_e\) its inclusion. On the trivial factor this is the section sum and common fibre relation. In a noncentral fibre the other two coefficient directions are acyclic; at the central leaves with no invariant vector the boundary section maps in the punctured base are injective by the surface pair sequence, and only the fibre kernel remains after the genuine fills. For clarity the resulting defect signatures are
\[
 \tau_E=\sum_v(\dim I_v)\sgn K_v-\sum_e\dim W_e-\sig Q_I,\qquad
 \tau_o=\sum_v\sgn K_v-|E|-\sig Q.
\]
Indeed on an anti-swap torus kernel use \(q_e=(y_{ve}-\beta_{ve}x_{ve})/\sqrt{\beta_{ve}}\); the other side is its transported signed image after the reciprocal scaling. With incidence \(Aq\) formed by summing these images on \(I_v\) with weights \(\sqrt\beta\), decomposition into the cap and component kernels gives \(K_I c=-Aq\). The Maslov form reduces to \(A^*K_I^{-1}A-I\); its complementary Schur form is \(K_I-AA^*=Q_I\). Ordinary coefficients have the same calculation with scalar invariant spaces. The anti-swap, rather than the ordinary graph of the seam isometry, is the kernel here.

In this example \(K=(4,-3/2,-3/2)\), each \(W_e\) is a line, and the only local invariant is the centre axis. Hence \(\tau_E=-1,\tau_o=-2\), giving defect \(\tau_E-3\tau_o=5\). Each filled leaf contributes \(-3\) by \eqref{eq:deletion}, since its local \(I\) vanishes and the new order-two lens has zero rho. On the filled centre compare with the bare zero-Euler triple (ordinary section fills at its two spine ports). Its adjoint plane is a conjugate pair of flat lines. In this comparison those line systems on the ports are acyclic; the two auxiliary opposite-orientation order-two lens line values vanish. The ordinary Maslov form is \(\sum_{\text{ports}} a_e b_e q_e^2\) on \(\sum b_e q_e=0\): subtract the ordinary section direction in each decomposition and use the zero-Euler sum of fibre components. This is positive of dimension one, so the filled central adjoint rho is twice its bare line rho minus two.

The bare line has holonomy \(e^{2\pi\ii t}\) on the fibre, and fractions \((u_j+t)/p_j\) on the cone loops. At \(t=0\) its two values are \(0,-2\): by \eqref{eq:deletion}, the first system is trivial, and the second has two order-four residues \(2\) giving \(2L(4,2)=-4\) on the conjugate plane. These bare line values persist to \(t=1/2\). To check the possible jump at zero, the Euler-zero Seifert form \(\xi\) can be chosen invariant and closed with fibre period one (the homomorphism \(h\mapsto1,z_j\mapsto1/p_j\) gives its de Rham class). In the local product metric \(d_t=d_0-2\pi\ii t\,\xi\wedge\). The nonzero fibre weights are acyclic near zero and all integer Fourier weights are acyclic for \(0<t<1\) by the contraction with the circle generator. Horizontal zero-weight harmonic lines split by type, with mixed operator slopes \(-2\pi t\) on \((1,0)\) and \(+2\pi t\) on \((0,1)\), since \(*_2(dx+\ii dy)=-\ii(dx+\ii dy)\). For the trivial line the base sphere has no horizontal harmonic one-forms, and the scalar/vertical pair has opposite signs. For the second line the base orbifold weights \((1/2,1/2,0)\) give invariant holomorphic degree \(-1\) on the sphere, as does the conjugate line; their \(H^1\)'s are zero. Thus no eta jump at zero in either case. The local eta variation on the flat cylinder vanishes (the only curvature is proportional to \(dt\wedge\xi\), with zero square). The interior values on the full star are consequently \(2(0)-2-6+5=-3\) and \(2(-2)-2-6+5=-7\). Along the chosen flat paths the odd-signature kernel dimension is constant on \(0<t<1\): the middle block has \(b_1=2\) on its common axis and each leaf \(b_1=5\) (five holed base of rank three, no invariant, then four line cap kills). The seam \(H^1\) map has rank three (two independent leaf section images plus one middle section direction); the \(H^0\) map has rank one from two seam lines. Global \(b_1=10\), \(b_0=0\). Thus no nonzero-mode crossing changes these absolute values along the paths.

\subsection{Exact lattice comparison and the wall limit}

The global Maslov exponent in \eqref{eq:finite} is \(-8\), since the cap sign sum is \(9\) and \(\sig Q=-1\). With \(t=a/r\), one leaf contributes \(L_r(a)=\sum_b C_2(b)^4 S_{ab}/S_{1b}^3\). Formula \eqref{eq:rec} gives exactly
\[
 L_r(a)=-(-1)^{r+a+1}e^{-5\pi\ii/r}r^2\,l_r(t),\quad
 l_r(t)=\frac8{r^3}\sum_{\substack{d\ {\rm odd}\\d<r}}
 \frac{\cos^4(\pi d/2r)}{\sin^3(\pi d/r)}
 e^{-\pi\ii d^2/r}\sin(\pi td).
\]
Thus \(l_r\to t(1-t)\) by the odd cubic sine series, uniformly on \([0,1]\). Indeed the lower weights are bounded by \(C/d^3\) and the upper range gives \(O(r^{-2})\), using the fourth-order cosine zero. Likewise \(l_r\) has uniformly bounded first derivative (multiply the weights by \(d\)), and for every fixed endpoint lattice defect \(f\), \(r l_r(f/r)\to f\) and \(r l_r(1-f/r)\to f\) (differentiate the dominated sine series at the endpoints, using odd \(d\)). In particular \(|l_r|\le C\min(t,1-t)\). The two leaf signs square out.

The remaining normalized centre term is
\[
 \frac{e^{-\pi\ii/4}}{4\sqrt2\,r}
 \sum_{a=1}^{r-1}\frac{l_r(t)^2}{\sin^3(\pi t)}
 \sum_{\substack{j,j'\bmod4\\\ell\bmod2}}
  P_{(2j,2j',2\ell)}(t)\,
     e^{-\pi\ii r(j^2+j'^2-2\ell^2)/2}
     e^{-\pi\ii a(j+j'-2\ell)/2},
\]
times a common \(1+O(1/r)\) from the finite twists. The resonant residues satisfy \(j+j'=2\ell\bmod4\). The four terms at each action \(0,1/2\) sum to \(\sin\pi t\), with the latter group multiplied by \((-1)^r\). For example, if \(x=\pi t/4\), the complementary residues each give \(-\sin^2 x\sin2x\) and \(+\cos^2 x\sin2x\) in pairs. The resonant Riemann sum therefore gives the bulk part, using \(\int_0^1[t(1-t)/\sin\pi t]^2\,dt=6\zeta(3)/\pi^4\). For the integral, integrate \(p(t)\csc^2(\pi t)\), \(p=t^2(1-t)^2\), twice by parts against \(\cot\pi t\) and \(\log\sin\pi t\); then \(p''=12(t^2-t+1/6)\) and the two cosine series give this value. Endpoint logarithms are square integrable, so the Fourier pairing is legitimate.

The nonresonant residues contribute the endpoint regimes. Uniformly, their summand with the oscillating phase removed is \(O(1/f)\) and its one-step difference is \(O(1/f^2)\), where \(f=\min(a,r-a)\), by the bounded derivative of \(l_r\) and its linear endpoint bound. The frequency gap is at least a quarter-period, so discrete parts make the tail with \(f>N\) uniformly \(O(1/N)\). Thus one may evaluate fixed defects first and then remove their cutoff. The positive endpoint is zero because its \(-2\) cap sine vanishes. At the negative endpoint each of four ordered angle pairs from \((\pi/4,3\pi/4)\), with third angle \(\pi/2\), lands at a pair-product interval wall. Pairing the two Weyl signs of each actual endpoint cap gives the sine sum \(\sum_{d\ge1}\sin(\alpha d)\sin(\beta d)\sin(\pi d/2)/d=\pi/8\): the three-sine Dirichlet sum is \(\pi/4\) in the strict pair interval and takes half its value at the wall. This elementary identity follows by product-to-sum and the linear sine series. The full magnitude of each tuple is \(1/(8\sqrt2\pi^2)\) with unit phase \(e^{5\pi\ii/4}\); explicitly the centre reciprocal constant \(e^{-\pi\ii/4}/(4\sqrt2)\) is multiplied by \((2\ii)^3\sin\alpha\sin\beta\sin(\pi/2)=-4\ii\), then by \(\pi^{-3}(\pi/8)\). Two endpoint tuples have phase one and two have the \(r\)-frequency \((-1)^r\), from the squares of their odd cap residues. The previously removed resonant sum vanishes pointwise at the wall, so this endpoint computation does not double count it. This proves \eqref{eq:star}.

The two specified branches have signed numerators
\[
 P_1=-\sin^2(x)\sin(2x),\qquad P_2=\cos^2(x)\sin(2x).
\]
Their unit quantum phases are \(e^{3\pi\ii/4},e^{-\pi\ii/4}\), exactly the two independently computed APS phases. At \(t=1\) they approach opposite simple pole terms with nonzero magnitudes; \(P_1+P_2=\frac12\sin\pi t\). For the same action and cap-class tuple their grouped density in \(dt=d(a/r)\) is
\[
 A_{\rm pair}(t)=\frac{e^{-\pi\ii/4}}{8\sqrt2}
                   \left[\frac{t(1-t)}{\sin\pi t}\right]^2 .
\]
Its boundary value is \(i\) times the central weight \(e^{5\pi\ii/4}/(8\sqrt2\pi^2)\) found above. This relates a limiting density to a coefficient in a separate charge regime, not an atom in that density. Neither a universal quarter-turn law nor the general absolute phase of every noncentral saddle is inferred.

\section{Marked comparison: a determinant and a loop term}\label{sec:marked}

For this section only, insert a smooth cutoff \(\chi(a/r)\) at each regular fibre, equal to one near the positive Weyl end and supported in a sufficiently small fixed neighborhood there. This is an actual finite Wilson-calculus insertion: the distinct fundamental eigenvalues \(2\cos(\pi a/r)\) allow polynomial interpolation of the cutoff values. It is not the unmarked invariant of Sections~\ref{sec:H}--\ref{sec:wall}. Use uniform coprime signed odd weights \(p,q\), both magnitudes at least three, and positive spine swaps; put \(D=-1/p-1/q\).

\begin{proposition}[Marked forest and loop benchmarks]\label{prop:marked}
For a connected tree with \(n\) vertices, \(Z_r^\chi/S_{11}\to (|pq|^n|\det Q|)^{-3/2}\). This is the inverse \(3/2\)-power of the absolute determinant of its full capped plumbing matrix. For the triangle (\(n=3\)), set \(t=(D^3-3D)/2\). The marked physical scalar instead has a nonzero order-one limit
\begin{equation}\label{eq:marked-tri}
 Z_r^\chi\longrightarrow -|pq|^{-9/2}
 \left[\frac1{2\sqrt{2(1-t)}}+\frac{\ii}{2\sqrt{2(1+t)}}\right].
\end{equation}
The two terms include an angular resonance and its principal value; the loop cannot be replaced by a three-copy tree determinant.
\end{proposition}

\begin{proof}
In the positive corner factor the linear part \(\pi a_va_w/r\) of each edge sine and retain its sinc. It cancels all inverse endpoint powers at each vertex. The noncentral reciprocal residue group is even on the signed lattice under simultaneous reversal of both cap residues and is zero at the origin (the full cap product and the removed central/central term each vanish there). Thus one-half the full signed Poisson sum gives its positive-grid sum, with no endpoint term. Each noncentral pair has nonzero linear frequency modulo two in units of \(\pi\) by the odd coprime hypothesis. Expanding incident sincs as averages of exponential linear steps, all other shifts are \(O(\delta)\) on support \(a/r<\delta\). Taking \(\delta\) sufficiently small, nonstationary parts eliminate these terms to arbitrary order, uniformly for all other charges. The remaining central pair contributes \(-\pi^2(a/r)^2/(pq)\) after its common reciprocal factor. Rescale \(a=\sqrt r\,x\); the compact-window Riemann sum gives the kernel
\[
 {\cal I}_G(D)=\int_{x>0}e^{\pi\ii D\sum_vx_v^2/2}
             \prod_v x_v^2\prod_{vw}\sinc(\pi x_vx_w)\,d^nx .
\]
Its integral is a smooth dyadic exhaustion, equivalent to Gaussian Abel damping. Here is the uniform tail test. In a box of largest charge scale \(B\gg\sqrt r\), separate the largest cluster of comparable scales by a sufficiently large fixed multiplicative gap. Expand its internal edge sines into signed exponentials; the true principal Hessian there is \(DI+\) signed adjacency, whose product by \(pq\) is \(-(p+q)I\bmod |pq|\) and hence invertible. Smaller neighboring charges contribute arbitrarily small relative shifts by choosing the gap large. A smooth subdivision selects a large coordinate with phase step \(\gtrsim B/r\), and \(N\) dyadic parts give \((r/B^2)^N\) against the worst normalized volume \(\prod_v(B_v/\sqrt r)^3\). Direct sine expansion on a large comparable scale costs only a fixed factor relative to its sinc envelope. Small boxes are summable by those cubic factors. This gives uniform lattice tails; the same integration by parts controls the continuum tails with or without positive Gaussian damping, whose dyadic seminorms are bounded uniformly. This justification precedes angular integration and does not presume integrability of its pointwise undamped determinant. Counting the lattice densities and common factors in \eqref{eq:finite} gives
\[
 r^{-3(l-n)/2}Z_r^\chi
 \ \longrightarrow\
 2^{n-l/2}\pi^{2n-l}|pq|^{-3n/2}e^{-3\pi\ii\sig Q/4}\,{\cal I}_G(D).
\]
The cap sign cancels the sign of \((-1/pq)^n\).

For the tree, write each sinc as the average of \(e^{\pi\ii x_vx_w n_v\cdot n_w}\) in unit directions; successively integrating leaf directions makes the joint spherical average exactly the product of sincs. With damping the Cartesian Gaussian is therefore \((4\pi)^{-n}\) times the integral with three copies of \(DI+\operatorname{Adj}G\). These matrices are invertible by the same modulus test. Their signature equals \(\sig Q\) by bipartite sign gauge, so their phase cancels the physical Maslov phase. The remaining prefactor is \(\sqrt2\pi r^{-3/2}(|pq|^n|\det Q|)^{-3/2}\); dividing by \(S_{11}\) gives the first assertion.

For the triangle, instead use independent Haar rotations on the edges inside the damped Gaussian; this gives every sinc without the leaf factorization. Gauge two edges along a tree to the identity; the third carries a single Haar \(SO(3)\) loop angle \(\theta\in[0,\pi]\) of density \((1-\cos\theta)/\pi\). This density follows from the unit-quaternion cover with scalar modulus \(\cos(\theta/2)\). Along its axis the Gaussian is the ordinary scalar \(B_0=DI+\operatorname{Adj} C_3\), of signature \(-1\) and determinant \(2(1+t)\). In its transverse plane the two real copies form a complex twisted cycle, with determinant \(2(t(z)+\cos\theta)\), \(t(z)=(z^3-3z)/2\) and \(z=D+2\ii\epsilon/\pi\). This follows directly from the two oriented triangle permutations of the phase twist. The plane Gaussian contributes \(2^3\ii^3/[2(t(z)+\cos\theta)]\). The axis square roots are continued separately from positive damping, giving \(2^{3/2}e^{-\pi\ii/4}/\sqrt{2(1+t)}\). Evaluate the angle integral at \(\epsilon>0\), before taking any boundary limit:
\[
 \frac1{2\pi}\int_0^\pi\frac{1-\cos\theta}{t(z)+\cos\theta}\,d\theta
   =\frac12\left[-1+\frac{1+t(z)}{\sqrt{t(z)^2-1}}\right],
 \qquad \sqrt{T^2-1}\sim T\text{ at infinity}.
\]
The identity is the tangent-half-angle integral with the specified branch. Since \(|D|<1\), \(|t|<1\) and \(t'(D)<0\), its boundary value is \(\tfrac12[-1+\ii(1+t)/\sqrt{1-t^2}]\). The underlying plane determinant has a zero at one real angle, so pointwise substitution would not give an absolutely convergent integral. Combining this boundary value with the axis factor, \(\sig Q=1\), and the common normalization gives \eqref{eq:marked-tri}. This establishes a distinct marked loop term without assigning it to the unmarked scalar.
\end{proof}

\section{What remains distinct}

In the H case the classical cubic and the actual leading closed logarithm are related by a density derived from a finite surgery sum. This does not allocate the next-order coefficient among equal-action characters. On the opposite-cap graph the absolute APS identity is topological even where the one-loop integrand is singular; an ordinary angular average still requires a separate integrability argument. At the central-fibre wall the correct quantum operation is to sum signed residues before taking the limit. A positive measure made by taking branch magnitudes first gives a different answer. Finally, the loop determinant calculation uses stated Wilson insertions. It does not give the unmarked leading coefficient or replace the signed regrouping at a Weyl wall.

\begin{thebibliography}{20}
\bibitem{APS} M.\,F.~Atiyah, V.\,K.~Patodi, and I.\,M.~Singer,
\emph{Spectral asymmetry and Riemannian geometry. II},
Math.\ Proc.\ Cambridge Philos.\ Soc.\ \textbf{78} (1975), 405--432.
\bibitem{Auckly} D.~Auckly, \emph{Topological methods to compute Chern--Simons invariants},
Math.\ Proc.\ Cambridge Philos.\ Soc.\ \textbf{115} (1994), 229--251,
\href{https://www.math.ksu.edu/about/personal/dav/topmethCS.pdf}{author's copy}.
\bibitem{GilmerMasbaum} P.\,M.~Gilmer and G.~Masbaum,
\emph{Maslov index, Lagrangians, mapping class groups and TQFT},
Forum Math.\ \textbf{25} (2013), 1067--1106,
\href{https://arxiv.org/abs/0912.4706}{arXiv:0912.4706}.
\bibitem{Jeffrey} L.\,C.~Jeffrey, \emph{Chern--Simons--Witten invariants of lens spaces and torus bundles, and the semiclassical approximation},
Commun.\ Math.\ Phys.\ \textbf{147} (1992), 563--604,
\href{https://webhomes.maths.ed.ac.uk/~v1ranick/papers/jeffrey.pdf}{article}.
\bibitem{KKR} P.~Kirk, E.~Klassen, and D.~Ruberman,
\emph{Splitting the spectral flow and the Alexander matrix},
Comment.\ Math.\ Helv.\ \textbf{69} (1994), 375--416,
\href{https://webhomes.maths.ed.ac.uk/~v1ranick/papers/kirklarub.pdf}{article}.
\bibitem{Massuyeau} G.~Massuyeau,
\emph{An introduction to the abelian Reidemeister torsion of three-dimensional manifolds},
Ann.\ Math.\ Blaise Pascal \textbf{18} (2011), 61--140,
\href{https://doi.org/10.5802/ambp.294}{doi:10.5802/ambp.294}.
\bibitem{SWW} H.~Sun, S.~Wang, and J.~Wu,
\emph{Self-mapping degrees of torus bundles and torus semi-bundles},
\href{https://arxiv.org/abs/0810.1799}{arXiv:0810.1799}, Prop.\ 1.5.
\bibitem{Toffoli} E.~Toffoli,
\emph{The Atiyah--Patodi--Singer rho invariant and signatures of links},
Proc.\ Edinburgh Math.\ Soc.\ \textbf{65} (2022), 404--440,
\href{https://arxiv.org/abs/2011.00525}{arXiv:2011.00525}.
\bibitem{Turaev} V.\,G.~Turaev,
\emph{Quantum Invariants of Knots and 3-Manifolds}, de Gruyter, 1994.
\bibitem{Witten} E.~Witten, \emph{Quantum field theory and the Jones polynomial},
Commun.\ Math.\ Phys.\ \textbf{121} (1989), 351--399.
\bibitem{Paper42B} \emph{Wilson Phase and Rational Sewing in Turaev--Viro Theory},
Project manuscript no.~42B, \href{https://github.com/g3173028019-beep/ai-automated-quantum-gravity-research/blob/main/docs/PAPER-INDEX.md}{manuscript index}.
\end{thebibliography}
\end{document}
